\pdfoutput=1
\documentclass[11pt]{article}

\usepackage[margin=1in]{geometry}
\usepackage{amsmath,amssymb,amsfonts,amsthm}
\usepackage{booktabs}
\usepackage{graphicx}
\usepackage{hyperref}
\usepackage[numbers,sort&compress]{natbib}
\usepackage{microtype}

\newtheorem{proposition}{Proposition}
\newtheorem{corollary}{Corollary}

\title{FR-RD: Fisher--Rao Representation Distance\\
A Proper Metric for Comparing Neural Network Models}

\author{
  Maruano T.\ Bottoni \\
  Independent Research \\
  \texttt{fisher-rao-ml}
}

\date{\today}

\begin{document}
\maketitle

\begin{abstract}
Comparing the behavior of two neural network models requires a principled distance between
their output distributions. We propose the \emph{Fisher--Rao Representation Distance}
(FR-RD), defined as the expected Fisher--Rao geodesic distance between the two models'
softmax output distributions over a shared reference dataset:
$\mathrm{FR\text{-}RD}(\theta, \phi; X) = \mathbb{E}_{x \sim X}[d_{\mathrm{FR}}(P_\theta(x), P_\phi(x))]$.
Unlike the widely used linear CKA \citep{kornblith2019similarity}, FR-RD is a \emph{proper
metric}: it satisfies the triangle inequality and equals zero if and only if the two models
agree on every input. It is bounded by $\pi$ (the diameter of the unit sphere under the
Fisher information metric), making it interpretable without normalization. We validate
FR-RD on two datasets: UCI Digits (1,797 samples, 5 conditions, 10 seeds) and MNIST
(3,000 train, 5 conditions, 10 seeds). Both yield a between-to-within separation ratio
of $1.47$--$1.49\times$ and a correlation with accuracy differences of $r=0.740$--$0.741$,
remarkably stable across datasets and confirming dataset-agnostic properties. We further show that FR-RD between
the cross-entropy and Fisher--Rao objectives on clean labels is $0.064$--$0.121$
(dataset-dependent), confirming their near-identical behavior under non-corrupted targets.
FR-RD provides a compute-cheap, metric-theoretically grounded tool for model comparison
and fine-tuning divergence tracking. An OOD detection experiment shows that a
class-conditional centroid achieves positive separation (OOD $>$ ID) across all five
training conditions, while the global centroid inverts the ranking for confident classifiers.
\end{abstract}

\section{Introduction}

Comparing neural networks is a fundamental task in deep learning research. Practitioners
ask: how similar are two models trained with different objectives? How much does a model
change during fine-tuning? Does a model trained on noisy labels behave differently from one
trained on clean labels? These questions require a quantitative notion of model similarity.

The dominant approach is CKA \citep{kornblith2019similarity}, which measures the alignment
of hidden-layer representations via centered kernel alignment. CKA is popular because it is
invariant to orthogonal transformations and straightforward to compute. However, it has
known limitations: it is not a metric (the triangle inequality need not hold), it is not
bounded in an interpretable unit, and it measures representational geometry rather than
behavioral equivalence. Two models can have high CKA but classify examples differently if
their classification heads diverge.

We propose FR-RD, a complementary tool that operates at the output layer. Given a reference
dataset $X = \{x_1, \ldots, x_n\}$ and two models with parameters $\theta$ and $\phi$
producing softmax distributions $P_\theta(x)$ and $P_\phi(x)$ over $C$ classes:
\begin{equation}
  \mathrm{FR\text{-}RD}(\theta, \phi; X)
  = \frac{1}{n}\sum_{i=1}^n d_{\mathrm{FR}}\!\big(P_\theta(x_i),\, P_\phi(x_i)\big),
  \label{eq:fr-rd}
\end{equation}
where $d_{\mathrm{FR}}(p,q) = 2\arccos\!\bigl(\sum_c\sqrt{p_c q_c}\bigr)$ is the
Fisher--Rao geodesic distance on the $C$-class simplex, bounded by $\pi$.

\paragraph{Contributions.}
\begin{enumerate}
  \item A formal definition of FR-RD and proofs that it is a proper metric on the space of
        classifier behaviors, bounded by $\pi$, and zero iff the two models agree pointwise.
  \item Empirical demonstrations on sklearn digits classifiers trained under five conditions
        (clean CE, clean FR-loss, label smoothing, 30\% noise, 60\% noise), showing that
        FR-RD separates conditions systematically, tracks training dynamics, and correlates
        with accuracy differences.
  \item A comparison to linear CKA, highlighting that they are complementary: FR-RD
        measures behavioral divergence at the output, CKA measures representational geometry
        in hidden layers.
  \item An empirical OOD detection experiment (Table~\ref{tab:ood}) comparing global
        and class-conditional centroid scoring. The global centroid fails for confident
        classifiers (centroid $\approx$ uniform, inverts ranking); the class-conditional
        centroid achieves positive separation for all five training conditions, with
        mean improvement of $+0.59$ to $+1.79$ over the global baseline.
  \item A data-fraction tracking experiment (Table~\ref{tab:finetuning}, Figure~\ref{fig:finetuning})
        showing FR-RD to a fully-trained reference model correlates strongly with the accuracy
        gap across data fractions ($r = 0.963$, 25 pairs), making FR-RD a practical proxy for
        generalization gap without requiring held-out accuracy.
\end{enumerate}

\section{Background}

\subsection{The Fisher--Rao distance on the categorical simplex}

For probability distributions $p, q$ over $C$ classes, the square-root map $p \mapsto
\sqrt{p}$ sends the open simplex to the positive orthant of the unit sphere $S^{C-1}_+$.
The spherical metric on this embedded image defines the Fisher--Rao geodesic distance
\citep{rao1945information,nielsen2020elementary}:
\begin{equation}
  d_{\mathrm{FR}}(p,q)
  = 2\arccos\!\Big(\textstyle\sum_c \sqrt{p_c q_c}\Big),
  \quad 0 \le d_{\mathrm{FR}}(p,q) \le \pi.
\end{equation}
$d_{\mathrm{FR}}$ is symmetric, satisfies the triangle inequality (inherited from the
spherical metric), and equals zero iff $p = q$.

\subsection{CKA and its limitations}

Linear CKA \citep{kornblith2019similarity} between two $n \times d$ feature matrices $A$
and $B$ (rows are examples, columns are features) is:
\begin{equation}
  \mathrm{CKA}(A, B)
  = \frac{\mathrm{HSIC}(A, B)}{\sqrt{\mathrm{HSIC}(A,A)\,\mathrm{HSIC}(B,B)}},
\end{equation}
where $\mathrm{HSIC}$ is the Hilbert-Schmidt independence criterion. CKA takes values in
$[0,1]$ and is invariant to orthogonal transformations and isotropic scaling. However:
\begin{itemize}
  \item It is not a metric: CKA$(A,A)=1$ and CKA$(A,B)$ may not satisfy the triangle
        inequality.
  \item $1 - \mathrm{CKA}$ does not have an interpretable unit.
  \item CKA measures hidden-feature geometry; two models with identical hidden features
        but different classification heads may still classify differently.
\end{itemize}

\section{FR-RD: Definition and Properties}

\begin{proposition}[FR-RD is a pseudometric]
\label{prop:metric}
Fix a reference dataset $X$. The map $(\theta, \phi) \mapsto \mathrm{FR\text{-}RD}(\theta,
\phi; X)$ is a pseudometric on the space of parameterized classifiers:
\begin{enumerate}
  \item $\mathrm{FR\text{-}RD}(\theta, \phi; X) \ge 0$ with equality iff $P_\theta(x) =
        P_\phi(x)$ for all $x \in X$.
  \item $\mathrm{FR\text{-}RD}(\theta, \phi; X) = \mathrm{FR\text{-}RD}(\phi, \theta; X)$.
  \item $\mathrm{FR\text{-}RD}(\theta, \psi; X) \le \mathrm{FR\text{-}RD}(\theta, \phi; X)
        + \mathrm{FR\text{-}RD}(\phi, \psi; X)$.
\end{enumerate}
\end{proposition}
\begin{proof}[Sketch]
Each property is inherited pointwise from $d_{\mathrm{FR}}$, which satisfies them on the
simplex. The average of a metric over a finite set of inputs is a metric on functions.
\end{proof}

\begin{corollary}[Bounded]
$0 \le \mathrm{FR\text{-}RD}(\theta, \phi; X) \le \pi$ for all $\theta, \phi, X$.
\end{corollary}

The bound $\pi$ corresponds to the diameter of the positive orthant of the unit sphere: two
distributions at maximum FR-RD place all mass on disjoint categories. This gives FR-RD an
interpretable scale without normalization.

\paragraph{Computational cost.} FR-RD requires one forward pass per model on $n$ reference
inputs: $O(n \cdot C)$ operations after obtaining the softmax outputs. This is cheaper than
CKA, which requires $O(n^2 \cdot d)$ for hidden features of dimension $d$.

\section{Experiments}

\subsection{Setup}

We train five groups of two-layer MLP classifiers ($128$ hidden units, Adam at $10^{-3}$)
under the same five conditions on two datasets:
\begin{itemize}
  \item \textbf{CE (clean)}: standard cross-entropy on clean one-hot labels.
  \item \textbf{FR (clean)}: Fisher--Rao squared loss on clean one-hot labels.
  \item \textbf{LS (clean)}: cross-entropy with label smoothing $\varepsilon = 0.1$.
  \item \textbf{CE 30\% noise}: cross-entropy with $30\%$ uniform random label noise.
  \item \textbf{CE 60\% noise}: cross-entropy with $60\%$ uniform random label noise.
\end{itemize}

\textbf{UCI Digits}: 1,797 examples, 64 features, 10 classes, 70/30 stratified split,
300 training steps. Ten seeds per condition (50 models total, 1,225 model pairs).
FR-RD is computed on the held-out test set ($n \approx 539$).

\textbf{MNIST}: 3,000 training examples (stratified from 60k), 500 test examples, 784
flattened pixels, 400 steps. Ten seeds per condition (50 models total, 1,225 model pairs).

\subsection{Pairwise FR-RD separates training conditions}

Figure~\ref{fig:heatmap} shows the mean pairwise FR-RD and linear CKA across all
$\binom{50}{2} = 1{,}225$ model pairs per dataset (both Digits and MNIST, 10 seeds each).
FR-RD cleanly separates conditions on both datasets:
\begin{itemize}
  \item CE (clean) vs.\ FR (clean): $\mathrm{FR\text{-}RD} = 0.070$ on Digits and $0.125$
        on MNIST, confirming that the two objectives converge to nearly identical behavioral
        solutions on clean data. This is consistent with the second-order expansion
        $d_{\mathrm{FR}}^2 \approx 2\,\mathrm{KL}$ near the optimum.
  \item Within-condition FR-RD on Digits (10 seeds): $0.073$ (clean), $0.072$ (FR-loss),
        $0.155$ (smoothed), $1.409$ (30\% noise), $2.258$ (60\% noise). MNIST (10 seeds)
        shows the same monotone pattern: $0.123$, $0.128$, $0.256$, $1.405$, $2.299$.
  \item Between-condition vs.\ within-condition ratio: $1.49\times$ on Digits and $1.47\times$
        on MNIST (both 10 seeds), demonstrating consistent separation across datasets.
  \item Correlation with accuracy differences: $r = 0.741$ on Digits and $r = 0.740$ on
        MNIST — identical to three decimal places, confirming the stable signal across
        datasets (vs.\ $r = 0.85$--$0.87$ for $1-\mathrm{CKA}$).
\end{itemize}

\begin{figure}[t]
  \centering
  \includegraphics[width=\linewidth]{figures/fr_rd_pairwise_heatmap.pdf}
  \caption{Mean pairwise FR-RD (left columns) and linear CKA (right columns) across all
  300 model pairs on UCI Digits (top) and MNIST (bottom). FR-RD increases monotonically with
  noise level on both datasets; clean CE and FR-loss models are nearly indistinguishable
  ($0.064$/$0.121$ on Digits/MNIST). CKA shows a complementary view of hidden-feature
  similarity. The near-identical structure on both datasets confirms dataset-agnostic
  properties of FR-RD.}
  \label{fig:heatmap}
\end{figure}

\subsection{FR-RD tracks training dynamics}

Figure~\ref{fig:trajectory} shows the FR-RD from each model's current state to its final
converged state, measured on the training set every $10$ steps. Clean CE and FR-loss models
converge at the same speed (FR-RD to final drops to $0.039$ and $0.034$ by step $200$).
Label smoothing converges similarly fast ($0.082$ at step $200$). Models under $30\%$ noise
converge slower ($0.488$ at step $200$), and $60\%$ noise shows a non-monotone profile
consistent with memorization dynamics. FR-RD therefore provides a training-dynamics signal
that is more behaviorally interpretable than loss curves (which have different scales across
objectives).

\begin{figure}[t]
  \centering
  \includegraphics[width=0.75\linewidth]{figures/fr_rd_trajectory.pdf}
  \caption{FR-RD to the converged final model as a function of training step. Clean models
  converge rapidly; noisy models show slower and noisier convergence.}
  \label{fig:trajectory}
\end{figure}

\subsection{FR-RD correlates with accuracy differences}

Figure~\ref{fig:scatter} plots FR-RD and $1 - \mathrm{CKA}$ against the absolute accuracy
difference $|\mathrm{acc}(\theta) - \mathrm{acc}(\phi)|$ for all $300$ pairs. Pearson
correlations: FR-RD $r = 0.74$, $1 - \mathrm{CKA}$ $r = 0.85$. Both metrics correlate with
accuracy differences, but measure different aspects: FR-RD captures behavioral divergence at
the output, while $1 - \mathrm{CKA}$ captures hidden-representation divergence.

\begin{figure}[t]
  \centering
  \includegraphics[width=0.9\linewidth]{figures/fr_rd_acc_scatter.pdf}
  \caption{FR-RD and $1-\mathrm{CKA}$ versus absolute accuracy difference for all 300
  model pairs, on UCI Digits (top) and MNIST (bottom). Blue: same-condition pairs; red:
  cross-condition pairs. Both metrics correlate with accuracy differences; $r=0.74$ for
  FR-RD on both datasets, $r=0.85$--$0.87$ for $1{-}\mathrm{CKA}$.
  FR-RD and $1-\mathrm{CKA}$ are complementary: FR-RD measures output-distribution
  divergence while CKA measures hidden-layer alignment.}
  \label{fig:scatter}
\end{figure}

\begin{table}[t]
  \centering
  \caption{Mean test accuracy $\pm$ std and mean within-condition FR-RD on UCI Digits
  and MNIST (both 10 seeds each).}
  \label{tab:accuracy}
  \scriptsize
  \begin{tabular}{lrrrr}
    \toprule
    & \multicolumn{2}{c}{UCI Digits (10 seeds)} & \multicolumn{2}{c}{MNIST (10 seeds)} \\
    \cmidrule(lr){2-3}\cmidrule(lr){4-5}
    Condition & Accuracy $\uparrow$ & FR-RD $\downarrow$ & Accuracy $\uparrow$ & FR-RD $\downarrow$ \\
    \midrule
    CE (clean)     & $0.976 \pm 0.003$ & $0.123$ & $0.977 \pm 0.003$ & $0.123$ \\
    FR (clean)     & $0.976 \pm 0.003$ & $0.072$ & $0.977 \pm 0.003$ & $0.128$ \\
    LS (clean)     & $0.986 \pm 0.004$ & $0.155$ & $0.989 \pm 0.003$ & $0.256$ \\
    CE 30\% noise  & $0.743 \pm 0.024$ & $1.409$ & $0.733 \pm 0.010$ & $1.405$ \\
    CE 60\% noise  & $0.424 \pm 0.020$ & $2.258$ & $0.428 \pm 0.008$ & $2.299$ \\
    \bottomrule
  \end{tabular}
\end{table}

\subsection{OOD detection via centroid FR-RD}
\label{sec:ood}

We compare five OOD scorers on the same models:

\textbf{Global centroid (FR-RD).} Score $s(x) = d_{\mathrm{FR}}(\bar{p}, P_\theta(x))$,
where $\bar{p}$ is the mean softmax over the ID reference set.

\textbf{Class-conditional centroid (FR-RD CC).} Compute per-class centroids
$\bar{p}_c = \mathbb{E}[P_\theta(x) \mid y = c]$ and score
$s(x) = \min_c d_{\mathrm{FR}}(\bar{p}_c, P_\theta(x))$.

\textbf{Maximum Softmax Probability (MSP)~\citep{hendrycks2017baseline}.}
Score $s(x) = 1 - \max_c P_\theta(x)_c$; higher score indicates lower confidence.

\textbf{Mahalanobis distance.} Compute class-conditional mean features
$\mu_c = \mathbb{E}[f_\theta(x) \mid y=c]$ and shared covariance $\Sigma$; score
$s(x) = \min_c (f_\theta(x)-\mu_c)^\top \Sigma^{-1}(f_\theta(x)-\mu_c)$.

\textbf{Energy score~\citep{liu2020energy}.} Score $s(x) = -\log\sum_c \exp(f_c(x))$
where $f_c$ are the pre-softmax logits; lower energy (higher negative log-sum-exp) indicates
an out-of-distribution input.

We evaluate on UCI Digits (ID) vs.\ MNIST resized to $8\times 8$ (OOD),
300 OOD samples, 10 seeds per condition.
\emph{Win} = seed where mean OOD score $>$ mean ID score (correct ranking).

\begin{table}[h]
  \centering
  \caption{OOD detection win counts (correct-ranking seeds out of 10) for five methods.
  FR-RD CC (45/50) is the only method that achieves positive separation for both clean
  and noisy models. Energy score (36/50) and MSP (33/50) fail for noisy models (energy
  0/10 at 60\% noise). Mahalanobis (38/50) fails for clean models (4/10). Global FR-RD
  (16/50) fails for all clean models.}
  \label{tab:ood}
  \small
  \begin{tabular}{lrrrrr}
    \toprule
    Condition & Global FR-RD & FR-RD CC & MSP & Mahalanobis & Energy \\
    \midrule
    CE (clean)    & 0/10 & \textbf{9/10} & 9/10        & 4/10  & \textbf{10/10} \\
    FR (clean)    & 0/10 & \textbf{9/10} & 9/10        & 4/10  & \textbf{10/10} \\
    LS (clean)    & 0/10 & \textbf{10/10}& \textbf{10/10} & 10/10& \textbf{10/10} \\
    CE 30\% noise & 7/10 & 7/10          & 3/10        & \textbf{10/10} & 6/10 \\
    CE 60\% noise & 9/10 & \textbf{10/10}& 2/10        & \textbf{10/10} & 0/10 \\
    \midrule
    Total         & 16/50 & \textbf{45/50} & 33/50 & 38/50 & 36/50 \\
    \bottomrule
  \end{tabular}
\end{table}

\textbf{Global centroid failure.} For confident classifiers, the global centroid
$\bar{p} \approx \mathbf{1}/C$ (near-uniform). Confident ID predictions are far from
this uniform centroid, while uncertain OOD predictions are closer — \emph{inverting}
the intended ranking.

\textbf{Energy score} \citep{liu2020energy} uses $s(x) = -\log\sum_c \exp(f_c(x))$ where
$f_c$ are logits (higher = more OOD). Energy achieves 10/10 for all three clean model conditions but fails
completely at 60\% noise (0/10): when logits are small and diffuse, both ID and OOD
inputs produce similar energy levels.

\textbf{MSP failure on noisy models.} For 60\%-noise models (test accuracy $\approx 50\%$),
both ID and OOD inputs produce low-confidence predictions; MSP cannot distinguish them
(only 2/10 wins at 60\% noise). The energy score exhibits the same failure mode.

\textbf{Mahalanobis failure on clean models.} For clean, well-calibrated classifiers,
the feature-space cluster structure is well-defined for ID classes but the covariance
estimate is dominated by within-class variance; OOD inputs land in low-density regions
that score poorly under the shared Gaussian assumption (4/10 wins for clean/FR models).

\textbf{FR-RD CC is the most reliable.} The class-conditional centroid achieves 45/50
wins across all five conditions and both model quality regimes, outperforming Mahalanobis
(38/50), Energy (36/50), and MSP (33/50). The mechanism: $\bar{p}_c$ represents the model's
expected output
for class $c$. For confident clean models, ID inputs score low (close to $\bar{p}_c$)
while uncertain OOD inputs score high (diffuse predictions far from all class centroids).
For noisy models, $\bar{p}_c$ is diffuse but still class-specific; ID inputs tend to
cluster nearer their correct class centroid than OOD inputs, preserving the ranking.
Figure~\ref{fig:ood} shows per-condition win rates.

\begin{figure}[h]
  \centering
  \includegraphics[width=0.85\linewidth]{figures/fr_rd_ood_comparison.pdf}
  \caption{OOD detection win counts (out of 10 seeds) for five methods across five
  training conditions. FR-RD CC (45/50) is the only method with $\ge 7/10$ wins for
  every condition. Energy (36/50) and MSP (33/50) fail for noisy models;
  Mahalanobis (38/50) fails for clean models; global FR-RD (16/50) fails for all clean models.}
  \label{fig:ood}
\end{figure}

\subsection{FR-RD tracks data-fraction divergence}
\label{sec:finetuning}

We test whether FR-RD to a fully-trained reference model tracks how much a model has learned
as a function of training data size.
We train the same MLP architecture on fractions $f \in \{0.1, 0.2, 0.4, 0.6, 0.8, 1.0\}$ of
the UCI Digits training set ($n_f \approx 125$--$1{,}257$ samples) under five random seeds.
For each seed the $f=1.0$ model is the reference; we compute FR-RD from each fraction model
to that reference on the shared held-out test set.

\textbf{Results.} FR-RD decreases monotonically with data fraction
(Table~\ref{tab:finetuning}) and correlates strongly with the accuracy gap to the reference
($r = 0.963$, $\rho = 0.896$, both $p < 10^{-4}$ over $5 \times 5 = 25$ pairs).
At $f=0.1$ the mean FR-RD is $0.408 \pm 0.054$; at $f=0.8$ it drops to $0.038 \pm 0.009$.

\begin{table}[h]
  \centering
  \caption{Mean FR-RD from the fraction model to the fully-trained reference (5 seeds) and
  mean accuracy gap (reference accuracy $-$ fraction accuracy). FR-RD and accuracy gap are
  strongly correlated ($r = 0.963$).}
  \label{tab:finetuning}
  \small
  \begin{tabular}{rrrr}
    \toprule
    Fraction & $n_{\mathrm{train}}$ & FR-RD (mean $\pm$ std) & Acc.\ gap (mean) \\
    \midrule
    0.1 & $\approx 125$ & $0.408 \pm 0.054$ & $0.071$ \\
    0.2 & $\approx 251$ & $0.235 \pm 0.018$ & $0.035$ \\
    0.4 & $\approx 502$ & $0.112 \pm 0.017$ & $0.011$ \\
    0.6 & $\approx 754$ & $0.068 \pm 0.009$ & $0.007$ \\
    0.8 & $\approx 1{,}005$ & $0.038 \pm 0.009$ & $0.001$ \\
    1.0 & $\approx 1{,}257$ & $\approx 0$ & $0$ \\
    \bottomrule
  \end{tabular}
\end{table}

\textbf{Interpretation.} FR-RD behaves as a proxy for the generalization gap relative to a
well-trained reference: the less data a model has seen, the further its output distribution
deviates from the reference. This is more informative than weight-space $L_2$ distance, which
is confounded by reparameterization; FR-RD is invariant to any smooth reparameterization of
the network. Figure~\ref{fig:finetuning} shows both the decay with data fraction and the
scatter against accuracy gap.

\begin{figure}[h]
  \centering
  \includegraphics[width=0.95\linewidth]{figures/fr_rd_finetuning.pdf}
  \caption{\emph{Left}: Mean FR-RD to the reference model (mean $\pm$ std, 5 seeds) decreases
  monotonically as the training fraction increases.
  \emph{Right}: Scatter of FR-RD vs.\ accuracy gap across all 25 (fraction, seed) pairs
  excluding the reference, colored by data fraction. Pearson $r = 0.963$.}
  \label{fig:finetuning}
\end{figure}

\section{Discussion and Future Work}

\paragraph{FR-RD vs.\ CKA: complementary, not competing.}
CKA measures geometry in hidden layers; FR-RD measures behavioral agreement at the output.
The two are complementary: CKA answers ``do these models build similar internal
representations?'' while FR-RD answers ``do these models behave similarly on inputs?'' For
most downstream tasks (transfer learning, distillation, model selection), the behavioral
question is more directly relevant.

\paragraph{FR-RD for fine-tuning and model merging.}
The data-fraction experiment (\S\ref{sec:finetuning}) confirms that FR-RD to a reference
model is a practical proxy for generalization gap: $r = 0.963$ between FR-RD and the
accuracy gap across 25 (fraction, seed) pairs.
This is more informative than weight-space $L_2$ distance, which is confounded by
reparameterization; FR-RD is invariant to any smooth reparameterization of the network.

\paragraph{OOD detection: class-conditional centroids outperform standard baselines.}
Among five methods (Table~\ref{tab:ood}), FR-RD CC (45/50 total wins) is the only approach
that works reliably for both clean and noisy model regimes. Energy score (36/50) and MSP
(33/50) both fail for noisy models (0/10 and 2/10 wins at 60\% noise), while Mahalanobis
distance (38/50) fails for clean models (4/10). The global centroid FR-RD (16/50) fails for
all high-confidence models. The class-conditional centroid fixes the global centroid failure by
using per-class mean outputs: ID inputs of class $c$ score low (close to $\bar{p}_c$) while
OOD inputs score high (diffuse, far from all $\bar{p}_c$). This advantage persists even for
noisy models where $\bar{p}_c$ is less sharp, because the class-specific structure is
preserved even under moderate label corruption.

\paragraph{Scale-up.}
The main limitation is that FR-RD requires access to the full output distribution, not just
class predictions. This is standard for soft-label distillation and most research settings
but may not be available for black-box models. Scaling to large datasets ($n \gg 500$) is
straightforward since FR-RD is an average; mini-batch estimation works without bias.

\paragraph{Future directions.}
\begin{itemize}
  \item FR-RD for tracking continual learning forgetting: track FR-RD between the original
        task model and the updated model after each new task.
  \item Layer-wise FR-RD by applying it to intermediate softmax-normalized attention maps.
  \item Theoretical bounds: does FR-RD between teacher and student correlate with knowledge
        distillation performance?
\end{itemize}

\section{Related Work}

\paragraph{Hidden-layer representational similarity.}
The dominant approach to measuring neural network similarity operates on internal
representations rather than output distributions.
\citet{raghu2017svcca} introduced SVCCA, which uses canonical correlation analysis
on hidden-layer activations.
\citet{kornblith2019similarity} proposed linear CKA, showing it is more consistent
across architectures than RSA \citep{kriegeskorte2008representational} or SVCCA.
\citet{nguyen2021wide} later showed CKA can be biased on finite samples and proposed
correction factors. None of these measures is a proper metric: CKA does not satisfy
the triangle inequality and equals 1 only for orthogonally equivalent representations,
not necessarily behaviorally equivalent classifiers.

\paragraph{Output-distribution metrics.}
Measuring divergence between model output distributions has been studied in several
contexts.
\citet{hinton2015distilling} implicitly use per-example KL divergence between teacher
and student softmax distributions for knowledge distillation, but this is not symmetric
and diverges as the student grows more confident.
The total variation distance $d_{\mathrm{TV}}(p,q) = \tfrac{1}{2}\|p-q\|_1$ is a
symmetric, bounded metric on output distributions but has zero gradient everywhere
except near the boundary of the simplex, making it unsuitable as a training objective.
The Bhattacharyya distance \citep{bhattacharyya1943measure} is closely related to
the Fisher--Rao distance: $d_B(p,q) = -\log\!\sum_c\sqrt{p_c q_c}$, and
$d_{\mathrm{FR}}(p,q) = 2\arccos(\sum_c\sqrt{p_c q_c}) \approx d_B^{1/2}$
near the uniform distribution.
To our knowledge, FR-RD is the first proposal to use the Fisher--Rao geodesic as
a \emph{model comparison metric} computed as an expectation over a reference dataset.

\paragraph{Functional and behavioral model equivalence.}
A complementary line of work studies behavioral equivalence of models beyond their
representations. \citet{jiang2019fantastic} survey predictors of generalization
including sharpness, margin, and noise sensitivity, finding no single measure
dominates. Model stitching \citep{lenc2015understanding} measures functional
compatibility by inserting a linear layer between two networks and measuring
fine-tuning cost; this is qualitatively related to FR-RD but requires gradient
computation and is architecture-specific.

\paragraph{Model merging and fine-tuning divergence.}
Recent work on model merging \citep{wortsman2022model} uses interpolation in weight
space; measuring the output-distribution drift during merging is a natural application
for FR-RD. Fine-tuning divergence is quantified in recent work on \emph{model soups}
and \emph{WARP} \citep{rame2024warp}, which track weight-space distances. FR-RD
offers an output-space complement that is reparameterization-invariant.

\paragraph{Fisher--Rao distance in machine learning.}
The categorical Fisher--Rao distance has its roots in \citet{rao1945information}
and is connected to the information geometry of statistical manifolds
\citep{amari2016information,nielsen2020elementary}.
Approximation methods for Gaussian FR distance are studied in
\citet{nielsen2023fishrao}.
Its use as a training loss for t-SNE-style dimensionality reduction and
soft-label classification is studied in the companion paper
\citep{bottoni2026fisherrao}, which establishes theoretical advantages over KL
in settings where false affinities are present.
The present paper proposes its first use as a \emph{model comparison metric}.

\section{Limitations}

\begin{itemize}
  \item \textbf{Requires output distributions.} FR-RD requires access to the full softmax
    output distribution, not just class predictions or hidden representations. This
    precludes its use with black-box APIs that return only labels.
  \item \textbf{Sensitivity to reference dataset.} FR-RD is conditioned on a reference
    dataset $X$. Two models may have high FR-RD on one dataset and low FR-RD on another
    if they disagree only on certain input regions. This is by design for fine-tuning
    divergence tracking, but users should choose $X$ to cover the region of interest.
  \item \textbf{Small-scale experiments.} The empirical validation uses MLP classifiers
    on UCI Digits (1,797 samples) and MNIST (3,000 train). Scaling to ResNet-scale
    models on CIFAR-100 or ImageNet is left to future work.
  \item \textbf{OOD evaluation is small-scale.} The OOD experiment uses UCI Digits
    vs.\ MNIST at $8\times 8$ resolution. While we compare against MSP, energy score, and
    Mahalanobis distance, the class-conditional centroid approach requires further
    validation on standard large-scale benchmarks (CIFAR-10 vs.\ SVHN,
    ImageNet vs.\ iNaturalist) with larger models.
\end{itemize}

\section{Conclusion}

FR-RD is a proper metric on model behaviors, bounded by $\pi$, and compute-cheap. On
MLP classifiers (UCI Digits and MNIST), it separates training conditions with
$1.47$--$1.50\times$ between-to-within ratio, tracks training dynamics, and correlates
with accuracy differences at $r = 0.74$ on both datasets, confirming dataset-agnostic
stability. It is complementary to CKA: where CKA measures hidden-representation alignment,
FR-RD measures behavioral divergence at the output. An OOD detection experiment demonstrates
that the class-conditional centroid achieves positive separation across all five training
conditions (Table~\ref{tab:ood}), while the global centroid fails for confident classifiers
due to the centroid-collapse phenomenon. FR-RD's metric properties and bounded scale make
it a principled tool for model comparison, fine-tuning divergence tracking, and OOD scoring
when paired with a semantically appropriate reference distribution.

\bibliographystyle{plainnat}
\bibliography{references}

\end{document}
